\begin{abstract}
Reliable perception in open-world driving requires both accurate dense
segmentation of known categories and the ability to flag unknown,
out-of-distribution (OoD) objects that were never seen during training.
This report studies road-scene understanding through the lens of modern
\emph{mask-classification} architectures, taking the Encoder-only Mask
Transformer (EoMT) with a DINOv2 backbone as the central model. We (i)
compare two provided EoMT checkpoints (one trained on Cityscapes for
semantic segmentation and one on COCO for panoptic segmentation) under a
single, fair 16-class semantic protocol obtained by projecting COCO
predictions into the Cityscapes label space; (ii) fine-tune the COCO model
on Cityscapes with a staged freezing schedule, recovering a validation mIoU
of $76,6\%$ with the backbone frozen except last two layers; and (iii) evaluate post-hoc anomaly
segmentation with Maximum Softmax Probability, Max-Logit, Max-Entropy and,
for the mask architecture, Rejected-by-All (RbA), across four public
benchmarks. The mask-based EoMT dramatically outperforms the pixel-based
ERFNet baseline, RbA gives the strongest and most balanced anomaly scores,
and fine-tuning the COCO model trades peak benchmark performance for markedly
more balanced behaviour across datasets. We discuss why these effects arise
and the limitations of the evaluation protocol.
\end{abstract}
